\documentclass[11pt,notitlepage]{report}
\usepackage[T1]{fontenc}
\usepackage[margin=1in]{geometry}
\usepackage{hyperref}
\hypersetup{colorlinks=true,linkcolor=blue,urlcolor=blue}
\title{Survey of Caching Strategies for Incremental Builds}
\author{FlashTeX Fixtures Team}
\date{}
\begin{document}
\maketitle
\begingroup
\let\clearpage\relax
\tableofcontents
\listoffigures
\listoftables
\endgroup
\chapter{Strategies and measurements}
\label{ch:strat}
Incremental build systems reuse work from previous compilations. The full measurements are in Section~\ref{sec:results} on page~\pageref{sec:results}; the cache layout is sketched in Figure~\ref{fig:cache} on page~\pageref{fig:cache}, and the raw timings are in Table~\ref{tab:times} on page~\pageref{tab:times}. Caching pays off only when dependency tracking is precise: over-approximate sets rebuild too much, while under-approximate sets produce stale output. This survey compares content hashing against timestamp comparison on workloads of increasing size.
\section{Background}
Content hashing skips unchanged files; timestamp comparison conservatively rebuilds them. The eviction policy of Section~\ref{sec:results} interacts with fingerprint granularity: fine-grained fingerprints admit more sharing but cost more to verify.
\begin{figure}[ht]
\centering
\rule{0.6\textwidth}{3cm}
\caption{Layout of the incremental build cache}
\label{fig:cache}
\end{figure}
\begin{table}[ht]
\centering
\begin{tabular}{|l|r|r|}
\hline
Workload & Files & Time (s) \\
\hline
Small & 120 & 1.2 \\
Medium & 1400 & 4.8 \\
Large & 12000 & 17.3 \\
\hline
\end{tabular}
\caption{From-scratch build times per workload}
\label{tab:times}
\end{table}
\section{Results}
\label{sec:results}
Content hashing wins on the large workload, cutting rebuild time roughly in half. As previewed at the start of this chapter, the gap grows with workload size. Figure~\ref{fig:cache} summarizes the winning layout and Table~\ref{tab:times} gives the baseline it improves upon.
\end{document}
